\newcommand{\developmenttext}{
\developmentIntroduction
\fourPrinciples
\openSourceCorrelation
\globalInfrastructure
\commonTechnologies
}

\newcommand{\developmentIntroduction}{As quatro liberdades do \textbf{Software Livre} ou modelos de desenvolvimento \textbf{Open Source} são princípios que garantem aos usuários controle e autonomia sobre os programas que utilizam. Essas liberdades permitem que qualquer pessoa utilize, estude, modifique, e compartilhe o software sem restrições impostas pelo desenvolvedor original. Elas são consideradas a base do conceito de software livre e incentivam a colaboração, o aprendizado e o desenvolvimento tecnológico. \par \fourPrinciples}

\newcommand{\fourPrinciples}{
A \textbf{Usar o Software}. Isso significa que o usuário pode utilizar o software e executa-lo da forma que desejar.Isso significa que o usuário pode utilizar o software da forma que desejar, seja para fins pessoais, acadêmicos, profissionais ou comerciais, sem precisar pedir autorização ao criador do programa. Essa liberdade garante que o software possa ser usado em diferentes contextos, sem limitações impostas pelo desenvolvedor.

A \textbf{Estudar o Software}, verificar como o programa funciona e adaptá-lo às próprias necessidades. Para que isso seja possível, é necessário que o \textbf{código-fonte} do software esteja disponível para os usuários. Dessa forma, programadores e estudantes podem analisar como o programa foi desenvolvido, entender sua estrutura e modificar partes do código para melhorar seu funcionamento ou adicionar novas funcionalidades.

A \textbf{Distribuir o Software}, redistribuir cópias do software permite que os usuários compartilhem o programa com outras pessoas, seja gratuitamente ou até mesmo cobrando por essa distribuição. Essa liberdade facilita o acesso à tecnologia e contribui para que mais pessoas possam utilizar e conhecer o software.

A \textbf{Modificar o Software},e disponibilizar essas modificações para outros usuários. Isso significa que, além de alterar o programa para uso próprio, o usuário também pode compartilhar versões modificadas com a comunidade. Dessa forma, melhorias feitas por um usuário podem beneficiar muitos outros, promovendo a evolução constante do software.}

\newcommand{\openSourceCorrelation}{ Atualmente, \textbf{Software Livre e Ferramentas Open Source} ocupam um papel central em diversas camadas da tecnologia da informação, estando presente desde sistemas operacionais até ferramentas de desenvolvimento e gerenciamento de projetos. Esse crescimento está diretamente relacionado aos seus princípios fundamentais, definidos pela Free Software Foundation e pela Open Source Initiative, que garantem a liberdade de uso, modificação e redistribuição do código, permitindo sua adaptação a diferentes contextos e necessidades \cite{fsf:2023, osi:2023}. Esse modelo colaborativo favorece a inovação contínua, contrastando com o modelo fechado de desenvolvimento tradicional \cite{raymond:2001}. Na prática, tecnologias amplamente utilizadas como Linux, MySQL e GitHub evidenciam como o software livre sustenta grande parte da infraestrutura digital moderna.

No contexto de ferramentas de gerenciamento de projetos, estudos demonstram que soluções open source, como OpenProject, ProjectLibre e Redmine, oferecem funcionalidades robustas como controle de tarefas, cronogramas, gestão de recursos e acompanhamento de progresso sendo comparáveis às ferramentas proprietárias como JIRA e Microsoft Project \cite{abramova:2016}. Ferramentas proprietárias tem uma tendência a apresentar interfaces mais intuitivas e melhor experiência de usuário, as soluções open source conseguem atender aos requisitos essenciais de gerenciamento de projetos, muitas vezes com maior flexibilidade e possibilidade de customização. Além disso o modelo open source também apresenta vantagens econômicas, reduzindo custos e incentivando a participação de comunidades no desenvolvimento e evolução dos sistemas\cite{lerner:2002}.

Por fim, observa-se que a adoção de software livre não se limita ao uso individual, mas se estende à infraestrutura global, sendo amplamente empregado em servidores, redes e dispositivos móveis, como no caso do sistema Android. Dessa forma, mesmo que usuários finais frequentemente optem por soluções proprietárias pela facilidade de uso, o software livre permanece como base tecnológica essencial, promovendo independência, transparência e inovação contínua no ecossistema digital.

}

\newcommand{\globalInfrastructure}{
O software de código aberto tornou-se um componente fundamental da infraestrutura digital moderna, sustentando a vasta maioria dos sistemas tecnológicos utilizados na indústria e nos serviços públicos. Estudos recentes realizados na Europa em Organizações de Setor Publico, indicam que mais de 90\% da infraestrutura digital global depende de componentes de código aberto, destacando seu papel crucial na viabilização de sistemas escaláveis, interoperáveis e colaborativos \cite{linaker:2026}. Além disso, o código aberto atua como o “tecido conjuntivo” da economia digital global, dando suporte a serviços essenciais como bancos, saúde e computação em nuvem \cite{forsgren:2021}. Essa ampla adoção reforça sua importância como base estratégica para ecossistemas tecnológicos resilientes e seguros \commonTechnologies.
}

\newcommand{\commonTechnologies}{ \par O papel central do software de código aberto na infraestrutura global se reflete diretamente na ampla adoção de tecnologias consagradas utilizadas no desenvolvimento de software moderno. Plataformas como o Android, baseado no kernel Linux, demonstram como componentes de código aberto impulsionam bilhões de dispositivos em todo o mundo, particularmente no ecossistema móvel. Da mesma forma, ferramentas como o Git\cite{git:2026} se tornaram o padrão para controle de versão distribuído, possibilitando o desenvolvimento colaborativo em escala entre organizações e comunidades abertas. Tecnologias de conteinerização como o Docker\cite{docker:2026} ilustram ainda mais esse impacto, fornecendo ambientes leves e portáteis que suportam aplicações escaláveis e nativas da nuvem. \par Mesmo em ambientes predominantemente proprietários, como o Microsoft Windows, há uma crescente integração com tecnologias de código aberto, incluindo suporte nativo para ferramentas e fluxos de trabalho de desenvolvimento baseados em Linux (WSL). Esses exemplos reforçam as conclusões que o software de código aberto não é apenas fundamental em um nível infraestrutural abstrato, mas também está concretamente incorporado às ferramentas, plataformas e ecossistemas cotidianos que sustentam o mundo moderno}